\section{Introducción: por qué la planeación es una capacidad central del coding agent}

Esta clase continúa la discusión sobre el diseño de herramientas para coding agents del tipo Codex / Claude Code. Quien da la clase señala que, al presentar antes el conjunto mínimo de herramientas de Codex, es fácil pasar por alto una herramienta que parece simple pero es muy importante: \texttt{update\_plan}. En una tarea de desarrollo real, un agent no enfrenta una sola pregunta con una sola respuesta, sino una tarea sobre un repositorio de código que se despliega de forma continua: primero entender el codebase, después formar un plan, luego ejecutar, verificar y ajustar, hasta terminar la tarea.

\begin{figure}[H]
\centering
\includegraphics[width=\textwidth]{frames/frame-00-00-50.jpg}
\caption{En la serie de Modern Agent, el Planning Agent y el diseño de herramientas de Codex forman un hilo continuo.\protect\footnotemark}
\end{figure}
\footnotetext{Intervalo de tiempo en el video: 00:00:45--00:00:55.}

\begin{importantbox}{Flujo típico de tareas de Codex}
En la abstracción de quien da la clase, Codex procesa una necesidad de desarrollo normalmente en tres pasos: \textbf{codebase exploration}, \textbf{plan} y \textbf{execution}. La exploration se encarga de leer el código y el contexto de forma intensiva, el plan divide el objetivo en pasos ejecutables, y la execution avanza de manera continua según la retroalimentación de la ejecución.
\end{importantbox}

Esta es también la razón por la que la planeación no es solo una función de la interfaz. En realidad es una representación de estado dentro del agent loop: el modelo necesita conservar en el contexto «qué hay que hacer ahora», «hasta qué paso se ha llegado» y «si hace falta replanear», y mostrarlo al usuario cuando sea necesario.

\subsection{Resumen de la sección}

Esta clase trata dos mecanismos complementarios de planeación en Codex: uno es el Plan Mode explícito, y el otro es la herramienta \texttt{update\_plan} durante la ejecución. El primero se orienta a la colaboración y a aclarar los requisitos; el segundo, a la gestión del progreso en tiempo de ejecución.
